\section{Variación de Parámetros en el Sistema de Primer Orden Asociado}

Considérese la ecuación diferencial lineal de segundo orden no homogénea:
\[
    u'' - a_1 u' - a_0 u = b(t) \label{eq:edo_original}
\] 
donde $a_0, a_1 \in \mathbb{C}$. Suponemos el caso de \textbf{raíz única (repetida)} para la ecuación característica, es decir, el discriminante es nulo: $a_1^2 + 4a_0 = 0$. La raíz es $\lambda = \frac{a_1}{2}$.

Definiendo el vector de estado $U = \begin{pmatrix} u \\ u' \end{pmatrix}$, la ecuación diferencial se escribe como el sistema:
\[
    U' = A U + B(t) \implies 
    \begin{pmatrix} u' \\ u'' \end{pmatrix} = 
    \begin{pmatrix} 0 & 1 \\ a_0 & a_1 \end{pmatrix} 
    \begin{pmatrix} u \\ u' \end{pmatrix} + 
    \begin{pmatrix} 0 \\ b(t) \end{pmatrix}
\]

Dado que $\lambda$ es una raíz con multiplicidad 2, las soluciones de la homogénea son $u_1 = e^{\lambda t}$ y $u_2 = t e^{\lambda t}$. La Matriz Wronskiana $W(t)$ se construye con estas soluciones y sus derivadas:
\[
    W(t) = \begin{pmatrix}
        e^{\lambda t} & t e^{\lambda t} \\
        \lambda e^{\lambda t} & e^{\lambda t} + \lambda t e^{\lambda t}
    \end{pmatrix}
\]
El determinante de esta matriz (el Wronskiano) es:
\[
    \det (W(t)) = e^{\lambda t}(e^{\lambda t} + \lambda t e^{\lambda t}) - (t e^{\lambda t})(\lambda e^{\lambda t}) = e^{2 \lambda t} \neq 0
\]

Buscamos una solución de la forma $U(t) = W(t)C(t)$, donde $C(t) = \begin{pmatrix} c_1(t) \\ c_2(t) \end{pmatrix}$. Sustituyendo en el sistema:
\[
    W'(t)C(t) + W(t)C'(t) = A W(t)C(t) + B(t)
\]
Como $W(t)$ es la M.F.S., cumple que $W'(t) = A W(t)$, por lo que el sistema se reduce a:
\[
    W(t)C'(t) = B(t) \implies 
    \begin{pmatrix}
        e^{\lambda t} & t e^{\lambda t} \\
        \lambda e^{\lambda t} & e^{\lambda t} + \lambda t e^{\lambda t}
    \end{pmatrix} 
    \begin{pmatrix} c_1'(t) \\ c_2'(t) \end{pmatrix} = 
    \begin{pmatrix} 0 \\ b(t) \end{pmatrix}
\]

Calculamos las derivadas de los parámetros:
\begin{itemize}
    \item Para $c_1'(t)$:
    \[
        c_1'(t) = \frac{1}{\det (W(t))} \begin{vmatrix} 0 & t e^{\lambda t} \\ b(t) & e^{\lambda t} + \lambda t e^{\lambda t} \end{vmatrix} = \frac{-b(t) t e^{\lambda t}}{e^{2\lambda t}} = -b(t) t e^{-\lambda t}
    \]
    \item Para $c_2'(t)$:
    \[
        c_2'(t) = \frac{1}{\det (W(t))} \begin{vmatrix} e^{\lambda t} & 0 \\ \lambda e^{\lambda t} & b(t) \end{vmatrix} = \frac{b(t) e^{\lambda t}}{e^{2\lambda t}} = b(t) e^{-\lambda t}
    \]
\end{itemize}

Integrando desde un punto inicial $t_0$, obtenemos:
\[
    c_1(t) = -\int_{t_0}^{t} b(s) s e^{-\lambda s} \, ds + x_1, \quad c_2(t) = \int_{t_0}^{t} b(s) e^{-\lambda s} \, ds + x_2
\]

La primera componente de $U(t) = W(t)C(t)$ nos da $u(t) = e^{\lambda t} c_1(t) + t e^{\lambda t} c_2(t)$. Sustituyendo los resultados:
\[
    u(t) = e^{\lambda t} \left( x_1 - \int_{t_0}^{t} b(s) s e^{-\lambda s} \, ds \right) + t e^{\lambda t} \left( x_2 + \int_{t_0}^{t} b(s) e^{-\lambda s} \, ds \right)
\]
Agrupando términos y simplificando la integral mediante la propiedad de linealidad:
\[
    u(t) = \underbrace{x_1 e^{\lambda t} + x_2 t e^{\lambda t}}_{\text{Sol. Homogénea}} + \underbrace{\int_{t_0}^{t} b(s) (t - s) e^{\lambda (t - s)} \, ds}_{\text{Sol. Particular}}
\]
